%%%%%%%%%%%%%%%%
%%% num 10 %%%
%%%%%%%%%%%%%%%%

\Chapter{10}

\Verse{1}
{
	\hP{וַיְדַבֵּר יְהֹוָה אֶל מֹשֶׁה לֵּאמֹר׃}
}
{
	\eP{
		10 And YHWH spoke to Moses, saying, [image] Numbers 10:2\\
		“Make two trumpets of silver. You shall make them hammered\\
		work. And you shall have them for calling the congregation\\
		and for having the camps travel. [image]\\
	}
}
{
}

\Verse{2}
{
	\hP{עֲשֵׂה לְךָ שְׁתֵּי חֲצוֹצְרֹת כֶּסֶף מִקְשָׁה תַּעֲשֶׂה אֹתָם וְהָיוּ לְךָ לְמִקְרָא הָעֵדָה וּלְמַסַּע אֶת הַמַּחֲנוֹת׃}
}
{
	\eP{}
}
{
}

\Verse{3}
{
	\hP{וְתָקְעוּ בָּהֵן וְנוֹעֲדוּ אֵלֶיךָ כּל הָעֵדָה אֶל פֶּתַח אֹהֶל מוֹעֵד׃}
}
{
	\eP{
		And when they will blow them, then all the congregation\\
		shall be gathered to you, to the entrance of the Tent of\\
		Meeting. [image]\\
	}
}
{
}

\Verse{4}
{
	\hP{וְאִם בְּאַחַת יִתְקָעוּ וְנוֹעֲדוּ אֵלֶיךָ הַנְּשִׂיאִים רָאשֵׁי אַלְפֵי יִשְׂרָאֵל׃}
}
{
	\eP{
		And if they will blow with one, then the chieftains, the\\
		heads of thousands of Israel shall be gathered to you.\\
		[image]\\
	}
}
{
}

\Verse{5}
{
	\hP{וּתְקַעְתֶּם תְּרוּעָה וְנָסְעוּ הַמַּחֲנוֹת הַחֹנִים קֵדְמָה׃}
}
{
	\eP{
		And when you will blow a blasting sound, then the camps\\
		that are camping on the east shall travel. [image]\\
	}
}
{
%	\footnote{
%		blasting. We do not know what sort of trumpet sound this was.
%	}
}

\Verse{6}
{
	\hP{וּתְקַעְתֶּם תְּרוּעָה שֵׁנִית וְנָסְעוּ הַמַּחֲנוֹת הַחֹנִים תֵּימָנָה תְּרוּעָה יִתְקְעוּ לְמַסְעֵיהֶם׃}
}
{
	\eP{
		And when you will blow a second blasting sound, then the\\
		camps that are camping on the west shall travel. They\\
		shall blow a blasting sound for their travels. [image]\\
	}
}
{
}

\Verse{7}
{
	\hP{וּבְהַקְהִיל אֶת הַקָּהָל תִּתְקְעוּ וְלֹא תָרִיעוּ׃}
}
{
	\eP{
		And for assembling the community you shall blow, but you\\
		shall not make a blast. [image]\\
	}
}
{
}

\Verse{8}
{
	\hP{וּבְנֵי אַהֲרֹן הַכֹּהֲנִים יִתְקְעוּ בַּחֲצֹצְרוֹת וְהָיוּ לָכֶם לְחֻקַּת עוֹלָם לְדֹרֹתֵיכֶם׃}
}
{
	\eP{
		And the sons of Aaron, the priests, shall blow the\\
		trumpets; and they shall be an eternal law for you through\\
		your generations. [image]\\
	}
}
{
}

\Verse{9}
{
	\hP{וְכִי תָבֹאוּ מִלְחָמָה בְּאַרְצְכֶם עַל הַצַּר הַצֹּרֵר אֶתְכֶם וַהֲרֵעֹתֶם בַּחֲצֹצְרֹת וְנִזְכַּרְתֶּם לִפְנֵי יְהֹוָה אֱלֹהֵיכֶם וְנוֹשַׁעְתֶּם מֵאֹיְבֵיכֶם׃}
}
{
	\eP{
		And when you will go to war in your land against a foe who\\
		afflicts you, then you shall blast the trumpets, and you\\
		will be brought to mind in front of YHWH, your God, and\\
		you will be saved from your enemies. [image]\\
	}
}
{
}

\Verse{10}
{
	\hP{וּבְיוֹם שִׂמְחַתְכֶם וּבְמוֹעֲדֵיכֶם וּבְרָאשֵׁי חדְשֵׁיכֶם וּתְקַעְתֶּם בַּחֲצֹצְרֹת עַל עֹלֹתֵיכֶם וְעַל זִבְחֵי שַׁלְמֵיכֶם וְהָיוּ לָכֶם לְזִכָּרוֹן לִפְנֵי אֱלֹהֵיכֶם אֲנִי יְהֹוָה אֱלֹהֵיכֶם׃}
}
{
	\eP{
		And on your happy occasion and at your appointed times and\\
		on your new moons, you shall blow the trumpets over your\\
		burnt offerings and over your peace-offering sacrifices,\\
		and they shall be for a commemoration for you in front of\\
		your God. I am YHWH, your God.” [image]\\
	}
}
{
}

\Verse{11}
{
	\hP{וַיְהִי בַּשָּׁנָה הַשֵּׁנִית בַּחֹדֶשׁ הַשֵּׁנִי בְּעֶשְׂרִים בַּחֹדֶשׁ נַעֲלָה הֶעָנָן מֵעַל מִשְׁכַּן הָעֵדֻת׃}
}
{
	\eP{
		And it was in the second year, in the second month, on the\\
		twentieth of the month that the cloud was lifted from over\\
		the Tabernacle of the Testimony. [image]\\
	}
}
{
}

\Verse{12}
{
	\hP{וַיִּסְעוּ בְנֵי יִשְׂרָאֵל לְמַסְעֵיהֶם מִמִּדְבַּר סִינָי וַיִּשְׁכֹּן הֶעָנָן בְּמִדְבַּר פָּארָן׃}
}
{
	\eP{
		And the children of Israel set out on their travels from\\
		the wilderness of Sinai. And the cloud tented in the\\
		wilderness of Paran. [image]\\
	}
}
{
%	\footnote{
%		the cloud tented. See the comment on 9:17.
%	}
}

\Verse{13}
{
	\hP{וַיִּסְעוּ בָּרִאשֹׁנָה עַל פִּי יְהֹוָה בְּיַד מֹשֶׁה׃}
}
{
	\eP{
		And they traveled from the first by YHWH’s word through\\
		Moses’ hand. [image]\\
	}
}
{
}

\Verse{14}
{
	\hP{וַיִּסַּע דֶּגֶל מַחֲנֵה בְנֵי יְהוּדָה בָּרִאשֹׁנָה לְצִבְאֹתָם וְעַל צְבָאוֹ נַחְשׁוֹן בֶּן עַמִּינָדָב׃}
}
{
	\eP{
		And the flag of the camp of the children of Judah traveled\\
		first by their armies, and over its army was Nahshon son\\
		of Amminadab. [image]\\
	}
}
{
}

\Verse{15}
{
	\hP{וְעַל צְבָא מַטֵּה בְּנֵי יִשָּׂשכָר נְתַנְאֵל בֶּן צוּעָר׃}
}
{
	\eP{
		And over the army of the tribe of the children of Issachar\\
		was Nethanel son of Zuar. [image]\\
	}
}
{
}

\Verse{16}
{
	\hP{וְעַל צְבָא מַטֵּה בְּנֵי זְבוּלֻן אֱלִיאָב בֶּן חֵלֹן׃}
}
{
	\eP{
		And over the army of the tribe of the children of Zebulun\\
		was Eliab son of Helon. [image]\\
	}
}
{
}

\Verse{17}
{
	\hP{וְהוּרַד הַמִּשְׁכָּן וְנָסְעוּ בְנֵי גֵרְשׁוֹן וּבְנֵי מְרָרִי נֹשְׂאֵי הַמִּשְׁכָּן׃}
}
{
	\eP{
		And the Tabernacle was taken down; and the children of\\
		Gershon and the children of Merari, who carried the\\
		Tabernacle, traveled. [image]\\
	}
}
{
}

\Verse{18}
{
	\hP{וְנָסַע דֶּגֶל מַחֲנֵה רְאוּבֵן לְצִבְאֹתָם וְעַל צְבָאוֹ אֱלִיצוּר בֶּן שְׁדֵיאוּר׃}
}
{
	\eP{
		And the flag of Reuben’s camp by their armies traveled,\\
		and over its army was Elizur son of Shedeur. [image]\\
	}
}
{
}

\Verse{19}
{
	\hP{וְעַל צְבָא מַטֵּה בְּנֵי שִׁמְעוֹן שְׁלֻמִיאֵל בֶּן צוּרִישַׁדָּי׃}
}
{
	\eP{
		And over the army of the tribe of the children of Simeon\\
		was Shelumiel son of Zurishadday. [image]\\
	}
}
{
}

\Verse{20}
{
	\hP{וְעַל צְבָא מַטֵּה בְנֵי גָד אֶלְיָסָף בֶּן דְּעוּאֵל׃}
}
{
	\eP{
		And over the army of the tribe of the children of Gad was\\
		Eliasaph son of Deuel. [image]\\
	}
}
{
}

\Verse{21}
{
	\hP{וְנָסְעוּ הַקְּהָתִים נֹשְׂאֵי הַמִּקְדָּשׁ וְהֵקִימוּ אֶת הַמִּשְׁכָּן עַד בֹּאָם׃}
}
{
	\eP{
		And the Kohathites, who carried the holy place, traveled;\\
		and they set up the Tabernacle by the time they came.\\
		[image]\\
	}
}
{
%	\footnote{
%		holy place. The Kohathites’ charge was the sacred objects that went
%		inside the holy place (the interior of the Tabernacle), namely: the ark,
%		the table, the menorah, and the altars and the equipment of the Holy
%		(Num 3:31). Footnote 10:21. they set up the Tabernacle by the time they
%		came. The Gershonite and Merarite priests traveled first and set up the
%		Tabernacle so that it was ready when the Kohathite priests arrived with
%		the ark, table, menorah, altars, and equipment.
%	}
}

\Verse{22}
{
	\hP{וְנָסַע דֶּגֶל מַחֲנֵה בְנֵי אֶפְרַיִם לְצִבְאֹתָם וְעַל צְבָאוֹ אֱלִישָׁמָע בֶּן עַמִּיהוּד׃}
}
{
	\eP{
		And the flag of the camp of the children of Ephraim by\\
		their armies traveled, and over its army was Elishama son\\
		of Ammihud. [image]\\
	}
}
{
}

\Verse{23}
{
	\hP{וְעַל צְבָא מַטֵּה בְּנֵי מְנַשֶּׁה גַּמְלִיאֵל בֶּן פְּדָהצוּר׃}
}
{
	\eP{
		And over the army of the tribe of the children of Manasseh\\
		was Gamaliel son of Pedahzur. [image]\\
	}
}
{
}

\Verse{24}
{
	\hP{וְעַל צְבָא מַטֵּה בְּנֵי בִנְיָמִן אֲבִידָן בֶּן גִּדְעוֹנִי׃}
}
{
	\eP{
		And over the army of the tribe of the children of Benjamin\\
		was Abidan son of Gideoni. [image]\\
	}
}
{
}

\Verse{25}
{
	\hP{וְנָסַע דֶּגֶל מַחֲנֵה בְנֵי דָן מְאַסֵּף לְכל הַמַּחֲנֹת לְצִבְאֹתָם וְעַל צְבָאוֹ אֲחִיעֶזֶר בֶּן עַמִּישַׁדָּי׃}
}
{
	\eP{
		And the flag of the camp of the children of Dan, the final\\
		one of all the camps, by their armies, traveled, and over\\
		its army was Ahiezer son of Ammishadday. [image]\\
	}
}
{
}

\Verse{26}
{
	\hP{וְעַל צְבָא מַטֵּה בְּנֵי אָשֵׁר פַּגְעִיאֵל בֶּן עכְרָן׃}
}
{
	\eP{
		And over the army of the tribe of the children of Asher\\
		was Pagiel son of Ochran. [image]\\
	}
}
{
}

\Verse{27}
{
	\hP{וְעַל צְבָא מַטֵּה בְּנֵי נַפְתָּלִי אֲחִירַע בֶּן עֵינָן׃}
}
{
	\eP{
		And over the army of the tribe of the children of Naphtali\\
		was Ahira son of Enan. [image]\\
	}
}
{
}

\Verse{28}
{
	\hP{אֵלֶּה מַסְעֵי בְנֵי יִשְׂרָאֵל לְצִבְאֹתָם וַיִּסָּעוּ׃}
}
{
	\eP{
		These were the orders of travel of the children of Israel\\
		by their armies when they traveled. [image]\\
	}
}
{
}

\Verse{29}
{
	\hJ{וַיֹּאמֶר מֹשֶׁה לְחֹבָב בֶּן רְעוּאֵל הַמִּדְיָנִי חֹתֵן מֹשֶׁה נֹסְעִים אֲנַחְנוּ אֶל הַמָּקוֹם אֲשֶׁר אָמַר יְהֹוָה אֹתוֹ אֶתֵּן לָכֶם לְכָה אִתָּנוּ וְהֵטַבְנוּ לָךְ כִּי יְהֹוָה דִּבֶּר טוֹב עַל יִשְׂרָאֵל׃}
}
{
	\eJ{
		And Moses said to Hobab, the son of Reuel the Midianite,\\
		Moses’ father-in-law, “We’re traveling to the place that\\
		YHWH said, ‘I’ll give it to you.’ Come with us, and we’ll\\
		be good to you, because YHWH has spoken of good regarding\\
		Israel.” [image]\\
	}
}
{
}

\Verse{30}
{
	\hJ{וַיֹּאמֶר אֵלָיו לֹא אֵלֵךְ כִּי אִם אֶל אַרְצִי וְאֶל מוֹלַדְתִּי אֵלֵךְ׃}
}
{
	\eJ{
		And he said to him, “I won’t go, but rather I’ll go to my\\
		land and to my birthplace.” [image]\\
	}
}
{
}

\Verse{31}
{
	\hJ{וַיֹּאמֶר אַל נָא תַּעֲזֹב אֹתָנוּ כִּי עַל כֵּן יָדַעְתָּ חֲנֹתֵנוּ בַּמִּדְבָּר וְהָיִיתָ לָּנוּ לְעֵינָיִם׃}
}
{
	\eJ{
		And he said, “Don’t leave us, because you know the way we\\
		should camp in the wilderness, and you’ll be eyes for us.\\
		[image]\\
	}
}
{
}

\Verse{32}
{
	\hJ{וְהָיָה כִּי תֵלֵךְ עִמָּנוּ וְהָיָה הַטּוֹב הַהוּא אֲשֶׁר יֵיטִיב יְהֹוָה עִמָּנוּ וְהֵטַבְנוּ לָךְ׃}
}
{
	\eJ{
		And it will be, when you go with us, that we’ll do good\\
		for you in proportion to the good that YHWH will do for\\
		us.” [image]\\
	}
}
{
}

\Verse{33}
{
	\hJ{וַיִּסְעוּ מֵהַר יְהֹוָה דֶּרֶךְ שְׁלֹשֶׁת יָמִים וַאֲרוֹן בְּרִית יְהֹוָה נֹסֵעַ לִפְנֵיהֶם דֶּרֶךְ שְׁלֹשֶׁת יָמִים לָתוּר לָהֶם מְנוּחָה׃}
}
{
	\eJ{
		And they traveled from the mountain of YHWH three days’\\
		journey. And the ark of the covenant of YHWH traveled in\\
		front of them three days’ journey to scout a resting place\\
		for them. [image]\\
	}
}
{
}

\Verse{34}
{
	\hJ{וַעֲנַן יְהֹוָה עֲלֵיהֶם יוֹמָם בְּנסְעָם מִן הַמַּחֲנֶה׃}
}
{
	\eJ{
		And YHWH’s cloud was over them by day as they traveled\\
		from the camp. [image]\\
	}
}
{
}

\Verse{35}
{
	\hJ{וַיְהִי בִּנְסֹעַ הָאָרֹן וַיֹּאמֶר מֹשֶׁה קוּמָה יְהֹוָה וְיָפֻצוּ אֹיְבֶיךָ וְיָנֻסוּ מְשַׂנְאֶיךָ מִפָּנֶיךָ׃}
}
{
	\eJ{
		And it was, when the ark traveled, that Moses said,\\
		“Arise, YHWH, and let your enemies be scattered, and let\\
		those who hate you flee from your presence,” [image]\\
	}
}
{
%	\footnote{
%		let your enemies be scattered. Moses’ ceremonial words that are
%		associated with the respective lifting and resting of the ark have been
%		taken as poetry in several scholarly treatments of Numbers: And it was,
%		when the ark traveled, that Moses said, “Arise, YHWH, and let your
%		enemies be scattered, and let those who hate you flee from your
%		presence,” Many identify this as poetry without any explanation. One
%		elaborates that it is “4:4 meter.” Given the complexities of speaking of
%		meter in biblical Hebrew poetry, and given that the lines in question
%		here do not add up to 4:4 by any metrical reckoning in any case, we are
%		better advised to be cautious of calling these lines poetry. Presumably
%		the element that has led some to assume that this is poetry is the
%		parallel of scattered enemies and fleeing adversaries, but simple
%		parallels within a single verse such as this may be found in both prose
%		and poetry in the Hebrew Bible. Some scholars have more circumspectly
%		spoken of these lines simply as “sayings.” The significance of this
%		point is that it directs us to recognize the effect of the late
%		introduction of short poems further on in the book (in Numbers 21–24).
%		More broadly, it should serve as a methodological caution in our
%		identification of verses as poetry in the Hebrew Bible. As a matter of
%		method, it is precarious to call a single verse in a prose context
%		poetry. Biblical poetry rarely rhymes and has no obvious meter. Its
%		best-known feature is the use of parallels (“Hear YHWH’s word, rulers of
%		Sodom. / Listen to our God’s teaching, people of Gomorrah,” Isa 1:10),
%		and it is also frequently characterized by puns and alliteration. But
%		parallels, puns, and alliteration are all to be found in biblical prose
%		as well. Without entering into a tedious discussion of the respective
%		definitions of prose and poetry in the Bible, we can say at this point
%		that in almost any passage that is more than a few verses long there
%		will be little scholarly disagreement over the simple identification of
%		the passage as prose or poetry. But making a judgment of a single verse
%		in the middle of a narrative and calling it poetry—and sometimes even
%		proceeding thereafter to analyze the author’s use of these poetic lines
%		in prose contexts as an artistic device—is a delicate and subjective
%		practice. This is not to say that there are no cases of single lines of
%		poetry embedded in biblical prose. It is rather a call for recognition,
%		as a matter of method, that whenever one identifies a verse as poetry
%		one must feel obliged to support that claim with analysis and evidence.
%		In the case of this passage, if I were to make such a case, it would be
%		based on the absence of prose particles (the definite article, the
%		direct object marker ’et, the relative pronoun ’ašer). But my point is
%		not to enter into this technical matter here, but rather to make the
%		larger point about poetry in the Torah in general.
%	}
}

\Verse{36}
{
	\hJ{וּבְנֻחֹה יֹאמַר שׁוּבָה יְהֹוָה רִבְבוֹת אַלְפֵי יִשְׂרָאֵל׃}
}
{
	\eJ{
		and when it rested, he said, “Come back, YHWH, the ten\\
		thousands of thousands of Israel.” [image]\\
	}
}
{
}
